\documentclass[11pt]{article}


\setlength{\parindent}{0pt}
\usepackage{xltxtra}
\usepackage{hyperref}
\hypersetup{hidelinks}
\usepackage{url}
\urlstyle{tt}
\usepackage{xcolor}
\definecolor{CVBlue}{RGB}{23,110,191}
\usepackage{calc}
\usepackage{graphicx}
\usepackage{tikz}
\usetikzlibrary{calc}
\usepackage{fontspec}
\usepackage{xeCJK}
\usepackage{enumitem}
\CJKsetecglue{} %% 取消中文与数字之间的间隙


%% 主文档字体设置
\setmainfont[
    Path = fonts/Main/,
    Extension = .otf,
    BoldFont = texgyretermes-bold.otf, % 加粗字体
]{texgyretermes-regular.otf} % 正文字体

% 中文字体设置
\setCJKmainfont[
    Path = fonts/hansans/,
    Extension = .ttf,
    BoldFont = NotoSansSC-Bold.ttf, % 加粗字体
]{NotoSansSC-Regular.otf} % 正文字体


\usepackage{relsize}
\usepackage{xspace}

% 使用 fontawesome（部分图标）
\usepackage{fontawesome} 

% A4纸，上下左右边距
\usepackage[
    a4paper,
    left=1.2cm,
    right=1.2cm,
    top=1.5cm,
    bottom=1cm,
    nohead
]{geometry}

\renewcommand{\baselinestretch}{1.5} % 行间距设为1.5

\usepackage{titlesec}
\usepackage{enumitem}
\setlist{noitemsep} % 取消列表项间的额外间距
%\setlist{nosep} % 取消所有垂直间距
\setlist[itemize]{topsep=0.25em, leftmargin=*}
\setlist[enumerate]{topsep=0.25em, leftmargin=*}

% --- 用于控制【不同项目之间】的垂直距离 ---
\newlength{\interProjectSpacing}
\setlength{\interProjectSpacing}{0.9em} % <--- 在此调整项目之间的距离
\newcommand{\projectsep}{\vspace{\interProjectSpacing}}

% --- 用于控制【项目标题】与下方【项目描述】的距离 ---
\newlength{\intraProjectTitleSep}
\setlength{\intraProjectTitleSep}{0.4em} % <--- 在此调整标题和描述的距离
\newcommand{\titlebreak}{\\[\intraProjectTitleSep]}

% --- 用于控制【项目描述】与下方【要点列表】的距离 ---
\newlength{\intraProjectListTopSep}
\setlength{\intraProjectListTopSep}{0.2em} % <--- 在此调整描述和列表的距离

% =======================================================================


\titleformat{\section}         % 定制 \section 命令 
{\large\bfseries\raggedright} % 将 section 标题设置为大号、粗体且左对齐
{}{0em}                      % 可用于为所有 section 添加前缀（如“章节...”）
{}                           % 可用于在标题前插入代码
[{\color{CVBlue}\titlerule}]  % 在标题后插入一条横线
\titlespacing*{\section}{0cm}{*1.6}{*1.2}



\begin{document}
\pagenumbering{gobble}

%%%% 利用tikz来定位照片，此处注释掉，可自行添加
% \begin{tikzpicture}[remember picture, overlay] 
%     \node[anchor = north east] at ($(current page.north east)+(-2cm,-1.2cm)$) {\includegraphics[height=3cm]{avatar.jpg}};
%   \end{tikzpicture}%
  %%%% 利用tikz来定位学校Logo，此处注释掉，可自行添加
  % \begin{tikzpicture}[remember picture, overlay] 
  %   \node[anchor = north west] at ($(current page.north west)+(0.5cm,+1.0cm)$) {\includegraphics[height=6cm]{zju.png}};
  % \end{tikzpicture}%
\centerline{\LARGE\bfseries{傅清祥}} 

\centerline{\normalsize{\faPhone\ 187xxxxx \quad \faEnvelopeO\ \href{mailto:qingxiang@xxx}{qingxiang@xxx}}} 

\centerline{\normalsize{\faGithubSquare\ \href{https://github.com/yourusername}{https://github.com/yourusername} \quad \faRssSquare\ \href{https://yourblog.github.io/}{https://yourblog.github.io}}} 
    
\section{\makebox[\widthof{\faGraduationCap}][c]{\color{CVBlue}\faGraduationCap}\ 教育背景}    
\textbf{浙江大学 物理学院} \hfill 2023.9 -- 至今\\[0.5em]
物理学（本科） \hfill 平均绩点：4.52/5 (90.59/100)
\begin{itemize}[nosep]
    \item 相关课程：计算物理 (5.0/100\%), 天体物理概论 (5.0/96\%), 热力学与统计物理 (4.8/94\%), 数学物理方法 (4.8/94\%)
    \item 英语水平：CET-4, CET-6 (567)
\end{itemize}

\section{\makebox[\widthof{\faUsers}][c]{\color{CVBlue}\faUsers}\ 科研经历}

% --- 第一个项目 ---
\textbf{行星系统早期共振演化与轨道稳定性分析} \hfill 2025.2 -- 2025.7 \titlebreak
项目描述：该项目聚焦于行星系统形成早期，因轨道迁移引发的平均运动共振（MMR）现象。通过数值模拟，研究共振捕获对行星轨道偏心率、倾角及系统长期稳定性的影响，旨在加深对系外行星系统构型多样性的理解。
\begin{itemize}[nosep, topsep=\intraProjectListTopSep]
    \item \textbf{数值模拟与计算}：独立完成了从 \textbf{Python} 到 \textbf{C++} 的语法迁移与工程化实践。基于 \textbf{REBOUND} N体积分框架，构建了多行星系统的轨道演化模型，并设计脚本实现了对共振角、轨道根数等关键参数的\textbf{自动化提取与批处理}。
    \item \textbf{数据分析与可视化}：运用 \textbf{NumPy} 和 \textbf{Matplotlib} 处理了数万条模拟时域数据，通过快速傅里叶变换（FFT）分析行星轨道的周期性变化。绘制了共振捕获概率随初始参数（如行星质量、粘滞盘密度）变化的\textbf{相图与热力图}，为后续理论分析提供了直观的数据支撑。
    \item \textbf{核心产出}：通过大量模拟，复现了经典的低质量行星易于捕获 2:1 共振的结论，并初步探索了共振锁定后系统稳定性的边界条件。
\end{itemize}

\projectsep

% --- 第二个项目 ---
\textbf{极端条件下应力样品托的优化设计与热-力耦合仿真} \hfill 2025.9 -- 2026.2 \titlebreak
项目描述：为满足课题组在低温高压极端环境下进行电输运测量的需求，对现有应力样品托的结构进行重新设计与力学性能评估，旨在提高样品在压力环境下的存活率与数据稳定性。
\begin{itemize}[nosep, topsep=\intraProjectListTopSep]
    \item \textbf{三维建模与设计优化}：自学并熟练运用 \textbf{SolidWorks} 进行三维实体建模。针对样品托的受力薄弱环节，提出了\textbf{倒角应力分散}与\textbf{材料局部加厚}的优化方案，并设计了新的定位夹具结构。
    \item \textbf{有限元仿真分析}：利用 \textbf{SolidWorks Simulation} 模块对设计模型进行了\textbf{静力学与热力学耦合仿真}。定量分析了样品托在不同材质（如铍铜、钛合金）与压力载荷下的形变位移与应力分布云图，从而筛选出最优结构参数与材料组合。
    \item \textbf{核心产出}：完成了从图纸设计到仿真验证的全流程，优化后的样品托结构仿真数据表明，其关键承压部位的应力集中现象减少了约 15\%。
\end{itemize}

\projectsep

% --- 第三个项目 ---
\textbf{基于图神经网络的脑网络拓扑结构与认知功能关联研究} \hfill 2026.4 -- 至今 \titlebreak
项目描述：本项目旨在利用深度学习技术，探索静息态功能核磁共振（rs-fMRI）数据构建的脑网络拓扑属性与个体认知能力之间的关系。尝试建立基于图神经网络的预测模型，以区分不同认知状态下的脑网络模式。
\begin{itemize}[nosep, topsep=\intraProjectListTopSep]
    \item \textbf{深度学习基础与复现}：系统自学了深度学习与神经网络基础理论，重点研读了图神经网络（GNN）在脑科学领域的应用文献。成功复现了 \textbf{SGN（Spatial Graph Network）} 与 \textbf{CNN（卷积神经网络）} 模型，用于处理脑网络的邻接矩阵与功能连接特征。
    \item \textbf{数据处理与建模}：参与课题组的脑影像数据预处理流程（基于 \textbf{DPABI} 与 \textbf{GRETNA} 工具包），提取了全局与局部的图论指标（如聚类系数、最短路径长度）。搭建了基于 \textbf{PyTorch} 的 GNN 框架，将脑网络结构转化为可学习的图数据，进行相关分类任务的初步训练与验证。
    \item \textbf{核心产出}：深入理解了图神经网络在非欧空间数据上的应用范式，积累了神经影像数据处理与深度学习建模的实战经验，为后续深入研究脑疾病生物标志物奠定了基础。
\end{itemize}

\section{\makebox[\widthof{\faCogs}][c]{\color{CVBlue}\faCogs}\ 技术栈}
\begin{itemize}[nosep]
    \item \textbf{编程语言：} Python, C++, Go
    \item \textbf{开发工具与框架：} VS Code, Git, SSH, Linux, PyTorch, REBOUND, SolidWorks
    \item \textbf{专业软件：} Matlab, Origin, DPABI, GRETNA
\end{itemize}

\section{\makebox[\widthof{\faGraduationCap}][c]{\color{CVBlue}\faList}\ 获奖情况}
\begin{itemize}
    \item 浙江大学校设一等奖学金（2024-2025学年，2025-2026学年）
    \item 学业优秀标兵、志愿服务标兵、文体活动标兵（2024-2025学年）
    \item 第xx届浙江省大学生田径锦标赛甲组女子5000m \textbf{第一名}
    \item 浙江省甲组女子5000m \textbf{纪录保持者}
\end{itemize}
    
\section{\makebox[\widthof{\faInfo}][c]{\color{CVBlue}\faInfo}\ 其他}
\begin{itemize}[parsep=0.5ex]
    \item \textbf{技术博客：} \href{https://yourblog.github.io/}{https://yourblog.github.io/}（记录一些物理与编程学习心得）
    \item \textbf{GitHub：} \href{https://github.com/yourusername}{https://github.com/yourusername}（包含部分科研项目代码与学习笔记）
    \item \textbf{兴趣爱好：} 长跑（国家一级运动员水平），阅读，开源社区贡献
\end{itemize}
\end{document}
